\section{Convergence Bounds for Stable Losses}

 \begin{theorem}\label{theo:RL_CV_bound}
For a generative flow trained on a target of known density with the stable $L^q$
flow-matching loss over a training probability measure that covers the background measure,
$\nuB\ll\nutrain$ with $\dln\in L^\infty(\nutrain)$, and run at inference with the corrected
terminal flow $\widehat\target := (\Finit+\Finstar-\Foutstar)^+$, the loss bounds two things
at once: the sampling error in total variation, and the defect between the total mass of the
initial flow and that of the target. That loss can be driven to zero over an $L^p$-universal
family with $q\leq p$ at an initial flow carrying the target's total mass, and over a family
whose initial mass is itself trained exactly when the target's total mass is a limit of the
masses that family realizes, under technical integrability assumptions.
\end{theorem}
The target's total mass is unknown in practice, so the mass clause is met by training rather
than imposed: an initial flow whose total mass is itself a parameter has that mass driven to
$\target(\Statespace)$ by the same loss that drives the residual --- the measure-theoretic
form of a GFlowNet learning the partition function $Z$ --- and the condition for it, that the
family realize masses approaching $\target(\Statespace)$, is necessary as well as sufficient.

The constant in the sampling-error bound is
$C=\frac{2}{\target(\Statespace)}\left\|\dln\right\|^{1/q}_{L^\infty(\nutrain)}\nuB(\Statespace)^{1-1/q}$,
and the mass defect carries its own; three things about $C$ are worth stating plainly. It
needs no normalization of $\target$, which may be an unnormalized reward as is standard. It
is scale-consistent: replacing
$\nuB$ by $c\nuB$ divides the loss by $c$ and multiplies $C$ by $c$, as it must, since
the left-hand side does not see $\nuB$. And it is \emph{not} free of the size of the state
space --- on a finite space with counting measure and uniform $\nutrain$ it reads
$2N/\target(\Statespace)$, so the effective bound carries a factor
$N^{1-1/q}$, which is a H\"older step and is attained when the residual is spread evenly.
At $q=1$ that factor is absent and the bound is dimension-free; Section \ref{sec:experiments}
measures both regimes. Full statement and proof: Theorem \ref{theo:RL_CV_bound_full},
Appendix \ref{app:stable}.

Two objects the practitioner chooses rather than inherits sit inside this bound: the
initial flow and the training measure. Appendix \ref{app:finit} works out the first --- the
certificate-optimal mass overshoots the partition function, undershoot is the dearer error,
and the optimal shape solves a one-sided Chebyshev problem --- and Appendix \ref{app:replay}
the second, where a replay buffer is read as a device shaping $\nutrain$ and the two pipelines
return opposite verdicts: uniform reweighting is optimal for the certificate's constant,
backward replay for the training rate.

\begin{theorem} \label{theo:IL_CV_bound_body}
For a probability initial flow, a normalized target and a forward policy whose star inflow
has a density, trained with the KL-weak-FM loss at any $q\in[1,+\infty]$ and an admissible
weight $b\geq1+\|\dln\|_{L^{q^*}(\nutrain)}$, the loss is bounded below by the self-entropy
$\mathcal H$ of the target --- with equality as soon as the family contains flows dominating
the target up to a factor $1-\varepsilon$ for every $\varepsilon>0$, attained at exact
matching, which strong $L^p$-universality supplies when the initial flow and the target have
densities in $L^p(\nuB)$ --- and the sampling error is
$O\bigl(\sqrt{\LossWKLFM^q-\mathcal H}\bigr)$ as the loss approaches $\mathcal H$.
\end{theorem}
This is the target-by-samples counterpart of Theorem \ref{theo:RL_CV_bound}: the loss never
uses the density of $\target$, only samples from it. Reaching the self-entropy may require a large $0$-flow component --- expected, since
reaching some regions of the state space can take arbitrarily many moves. Full
statement and proof: Theorem \ref{theo:IL_CV_bound}, Appendix \ref{app:stable}.
